\section{Shannon 熵与二元熵}
\label{sec:07-Information-Theory/01-Information-and-Entropy/02}

朋友每天给你发一条微信，内容只有一个字：``晴''或``雨''。你所在的城市一年里 99\% 的日子是晴天。今天收到``晴''，你眼皮都没抬——这消息值几个钱？明天万一收到``雨''，你可能要翻箱倒柜找伞。

但这个感觉该怎么变成数字？

\subsection{从惊讶到自信息}

首先解决一个更简单的问题：单次消息的``信息量''。你本能觉得，小概率事件发生了，消息就值钱；大概率事件，知道了也不新鲜。试一个候选公式：

\[
\imath_X(X)=-\log P(X).
\]

几个直觉检查：概率越小，$-\log$ 越大，对得上``罕见事件信息多''；若 $P=1$，必然事件的信息量为零——还没看消息就已经知道内容，确实零收货。对数底取二时，单位是 bit：一个公平硬币的正反面各带 $-\log_2(1/2)=1$ bit。

这里暂时不追问对数为什么要这样放——现在只需要确认：这个定义和我们日常说``这消息有多少料''的直觉一致。

\subsection{消息还没来，平均要准备多少信息}

自信息 $\imath_X(X)$ 本身还是随机变量：消息内容不确定，信息量也不确定。我们真正想知道的，是在收到消息\emph{之前}，平均每次能获得多少新信息。自然做法是对可能的消息取期望：

\[
H(X)
=\E[\imath_X(X)]
=-\sum_x p(x)\log p(x).
\]

这就是 Shannon 熵。约定 $0\log 0=0$（概率为零的事件对期望没有贡献）。对数底为二时，单位是 bit。

\subsection{两个极端帮你建立标尺}

若 $X$ 几乎必然为某个值，$H(X)=0$。消息还没发你就知道内容——平均获得零信息。极端确定对应零熵。

若 $X$ 在 $m$ 个结果上均匀分布，$H(X)=\log m$。所有结果等可能，没法优先猜任何一个——不确定性最大。此时熵恰好等于结果种类数的对数。

对任何固定 $m$ 个可能值的分布：

\[
0\le H(X)\le\log m,
\]

上界当且仅当分布均匀时达到。这个上界本身值得推导——将在最大熵一节用 Lagrange 乘子法严格验证。但可以先接受一个粗略的直观：均匀分布让人完全无偏，其他分布都带着``倾斜感''，自然缺了一些不确定性。

\subsection{偏置硬币：二元熵}

现在回到两结果情形。设 $X\sim\mathrm{Bernoulli}(p)$，熵是 $p$ 的函数：

\[
h_2(p)
=-p\log_2 p-(1-p)\log_2(1-p).
\]

这个函数叫二元熵函数。它关于 $p=1/2$ 对称：$h_2(p)=h_2(1-p)$。在 $p=0$ 或 $p=1$ 时熵为零（完全确定），在 $p=1/2$ 时达到最大值 1 bit。

一个容易被忽略的点：只要结果仍有两种，并不代表一定需要 1 bit 来编码每个符号。$p=0.1$ 时 $h_2(0.1)\approx 0.47$ bit——平均信息确实远小于 1 bit。这正是上一节预告的：偏置硬币的``平均信息可小于 1 bit''。单符号编码受整数位限制，但通过对长序列块联合编码，每位符号的平均开销可以逼近这个非整数熵。

所以别把``两种可能结果''和``需要 1 bit''划等号。两个结果只给出一个上界；实际需要多少 bit，由概率分布的倾斜程度决定。

\subsection{混合分布会让熵变大}

假设有两个分布 $p$ 和 $q$，随机选一个来采样：以概率 $\lambda$ 选 $p$，以 $1-\lambda$ 选 $q$，得到混合分布

\[
r=\lambda p+(1-\lambda)q,\qquad 0\le\lambda\le 1.
\]

那么

\[
H(r)\ge\lambda H(p)+(1-\lambda)H(q).
\]

不知道样本来自哪个成分，这层额外的不确定性体现为熵的增加。严格来说：每一项 $-p_i\log p_i$ 关于 $p_i$ 是凹函数（二阶导为负），求和不改变凹性，因此整体熵在概率单纯形上是凹的。

凹性还有一个优雅的重新解读。均匀分布在固定支撑上最大化熵：把任意分布的所有坐标随机置换再取平均，得到的就是均匀分布；而凹性保证均值熵不小于原分布熵，所以均匀分布熵最大。不需要 Lagrange 乘子也能看出结论，只是推导方向不同。

\subsection{把变量压扁不会增加熵}

若 $Y=g(X)$ 是 $X$ 的确定函数：

\[
H(Y)\le H(X).
\]

函数可能把多个不同 $x$ 映射成同一个 $y$——合并了可区分的结果，不创造新的随机选择。严格证明需要条件熵工具，留在后续单元。若 $g$ 在 $X$ 的支撑上一一对应，则熵保持不变：可逆变换只是重命名，不确定度没有实质变化。

一个易混淆之处：把离散变量编码成更长的二进制字符串，物理位数增加了，但随机变量的熵并不因此增加。熵关心的是可区分的随机结果及其概率，不关心表示方式的冗余度。一个四结果的均匀变量，无论用 $\{00,01,10,11\}$ 还是用 $\{\text{春},\text{夏},\text{秋},\text{冬}\}$ 表示，熵都是 $2$ bit。

\subsection{熵就是压缩的硬极限}

要落实``熵衡量信息''的直觉，最直接的办法是回到编码问题。前缀码要求任何码字都不是另一个码字的前缀——读到一半不需要回头猜测，可以即时解码。用二进制前缀码且熵以 bit 计时：

\[
H(X)\le L,
\]

$L$ 是平均码长。任何前缀码的平均长度不可能低于熵。

反过来，对有限字母表，Huffman 码等最优前缀码满足

\[
L<H(X)+1.
\]

所以熵不仅是从 $-p\log p$ 公式算出来的抽象数字——它是无失真压缩所需平均位数的精确下界，且总是可以通过对长序列块编码逼近。两个数字走到了一起，定义就有了物理支点。

\subsection{熵看不见分布的形状明细}

熵是分布的一个摘要数字。两个完全不同形状的分布可以有相同熵——熵不记录哪个结果叫``晴''、哪个叫``雨''，也不在乎标签语义。已知 $H(X)=1$，不足以重建 $p(x)$。

离散熵对可逆重命名不变；但连续变量的微分熵连缩放都会改变。这个差异引出下一节的核心问题。
